\chapter{Ambiente e comandi di riproduzione}
\label{app:riproduzione}

\section{Struttura del repository}

\begin{lstlisting}[style=shell, numbers=none]
quantum-ml-research/
|-- run.py                      # entry point comune ai quattro esperimenti
|-- pyproject.toml              # manifest unico delle dipendenze (uv)
|-- uv.lock                     # ambiente pinnato, autoritativo
|
|-- quantum_framework/          # libreria condivisa, installata come package
|   |-- layers/quantum_linear.py
|   |-- data/medmnist_loader.py
|   |-- training/trainer.py
|   |-- evaluation/{metrics.py, statistics.py}
|   `-- utils/{config.py, experiment.py}
|
|-- experiments/
|   |-- 01_quantum_gpt/         # proiezioni Q,K,V quantistiche
|   |-- 02_quantum_vit/         # patch embedding quantistico
|   |-- 03_quantum_reg/         # ablazione sulla regolarizzazione
|   `-- 04_quantum_kernel/      # kernel di fedelta vs similarita coseno
|
|-- tests/                      # test_config.py, test_metrics.py, test_statistics.py
`-- docs/
\end{lstlisting}

\section{Installazione}

\begin{lstlisting}[style=shell, numbers=none]
git clone <url-del-repository>
cd quantum-ml-research
uv sync               # dipendenze principali (esperimenti 02, 03, 04)
uv sync --extra gpt   # aggiunge torchinfo, torchviz, datasets, transformers (esp. 01)
uv sync --extra gpu   # aggiunge pennylane-lightning[gpu]
\end{lstlisting}

L'ambiente pinnato è \file{uv.lock} ed è la fonte autoritativa;
\file{requirements.txt} è un riassunto leggibile e non va usato per installare.
Per la simulazione accelerata su GPU, oltre all'extra \code{gpu}, occorre
impostare \code{q\_device = "lightning.qubit"} nella configurazione al posto di
\code{default.qubit}.

\section{Controlli preliminari}

Prima dell'esecuzione si controllano albero di lavoro e test:

\begin{lstlisting}[style=shell, numbers=none]
git status --porcelain                                  # deve essere vuoto
uv run python -m unittest discover -s tests -v
uv run pre-commit run --all-files                       # ruff: lint e formattazione
\end{lstlisting}

Prima di una run completa si esegue inoltre uno \emph{smoke test} con poche
epoche, controllando forme dei tensori, checkpoint e schema delle metriche.

\section{Invarianti di configurazione}
\label{app:invarianti}

Il metodo \code{\_\_post\_init\_\_} di \classe{BaseViTConfig}
(\cref{sec:base-config}) verifica $H \bmod p = 0$,
$\dvit \bmod \nhead = 0$ e la positività di dimensioni e profondità. In caso di
violazione solleva \code{ValueError} prima di costruire il modello.

\section{Valutazione di robustezza al rumore}
\label{ssec:rumore}

Il metodo \code{\_test\_noise\_robustness()} di \classe{BaseTrainer} valuta il
modello addestrato su due tipi di corruzione, rumore gaussiano e \emph{sale e
pepe}, a cinque intensità $\{0{,}05,\ 0{,}1,\ 0{,}2,\ 0{,}3,\ 0{,}5\}$, riportando
loss, accuracy, AUC e F1 per ogni combinazione. Il seed delle corruzioni è
derivato in modo deterministico da quello della run, secondo
\code{seed + 10000 + indice\_tipo·100 + indice\_intensità}.

Le varianti con lo stesso seed ricevono le stesse corruzioni. La run definitiva
dell'esperimento~03 non include questa valutazione.

\section{Esperimento 01 --- Quantum GPT}

\begin{lstlisting}[style=shell, numbers=none]
uv run python run.py 01 --mode compare --config thesis \
  --dataset data/input.txt \
  --seeds 42,137,256,512,1024,2048,4096,8192,16384,32768 \
  --name thesis_shakespeare --force_cpu --resume
\end{lstlisting}

Corpus, preset e seed devono restare invariati fra le due varianti. L'aggregato
è scritto in \file{experiments/comparison\_thesis\_shakespeare.json}; la copia
congelata è \file{docs/checkpoints/experiment\_01\_quantum\_gpt\_n10.json}.

\section{Esperimento 02 --- Quantum ViT}

\begin{lstlisting}[style=shell, numbers=none]
uv run python run.py 02 compare --dataset pathmnist --epochs 50 \
  --seeds 42 137 256 512 1024 --name thesis_pathmnist \
  --q-device default.qubit --force-cpu
\end{lstlisting}

In caso di interruzione si ripete lo stesso comando aggiungendo \code{--resume}:
il checkpoint atomico conserva modello, ottimizzatore, scheduler, storico, RNG e
stato del generatore del DataLoader. PathMNIST è il test primario; le eventuali
repliche su altri dataset sono analisi secondarie e mantengono invariati gli
iperparametri fra le varianti.

\section{Esperimento 03 --- Ablazione sulla regolarizzazione}

\begin{lstlisting}[style=shell, numbers=none]
uv run python run.py 03 ablation --epochs 400 --train-subset 100 \
  --seeds 42 137 256 512 1024 \
  --datasets pathmnist bloodmnist dermamnist \
  --dropout 0 --weight-decay 0 --eval-interval 40 --skip-noise
\end{lstlisting}

La run comprende $3 \times 5 \times 3 = 45$ run e salva
\file{ablation\_progress.json} dopo ciascuna. In caso di interruzione:

\begin{lstlisting}[style=shell, numbers=none]
uv run python run.py 03 ablation --epochs 400 --train-subset 100 \
  --dropout 0 --weight-decay 0 --eval-interval 40 --skip-noise \
  --resume experiments/ablation_thesis_regularization.progress.json
\end{lstlisting}

\section{Esperimento 04 --- Kernel quantistico}

\begin{lstlisting}[style=shell, numbers=none]
uv run python run.py 04 --seed 42 --repeats 1 \
  --n-class0 80 --n-class1 20 --n-qubits 8 \
  --feature-map iqp --iqp-repeats 2 --kernel-method statevector
uv run python run.py 04 --seed 42 --repeats 20 \
  --n-class0 80 --n-class1 20 --n-qubits 8 \
  --feature-map iqp --iqp-repeats 2 --kernel-method pairwise
\end{lstlisting}

Ogni ripetizione usa il seed $\code{seed} + r$ e salva gli indici selezionati e
la varianza spiegata dalla PCA, oltre alle matrici dei kernel. I controlli di
separabilità usano il metodo a vettore di stato:

\begin{lstlisting}[style=shell, numbers=none]
uv run python run.py 04 --seed 42 --repeats 20 \
  --n-class0 50 --n-class1 50 --n-qubits 8 \
  --feature-map iqp --iqp-repeats 2 --kernel-method statevector
uv run python run.py 04 --seed 42 --repeats 20 \
  --n-class0 27 --n-class1 73 --n-qubits 8 \
  --feature-map iqp --iqp-repeats 2 --kernel-method statevector
\end{lstlisting}

La validazione SVC è eseguita dalla directory
\file{experiments/04\_quantum\_kernel} nei regimi 80/20, 50/50 e 27/73:

\begin{lstlisting}[style=shell, numbers=none]
uv run python classifier_validation.py --n-class0 80 --n-class1 20 \
  --repeats 20 --seed 42
uv run python classifier_validation.py --n-class0 50 --n-class1 50 \
  --repeats 20 --seed 42
uv run python classifier_validation.py --n-class0 27 --n-class1 73 \
  --repeats 20 --seed 42
\end{lstlisting}

\section{Output di una run}

Ogni esecuzione scrive in una directory identificata da un timestamp:

\begin{itemize}
  \item \file{config.json} --- snapshot completo della configurazione, comprensivo
    delle proprietà calcolate;
  \item \file{run\_manifest.json} --- provenienza: commit, flag \emph{dirty},
    riga di comando, timestamp, versioni di librerie e hardware;
  \item \file{params.json} --- conteggio dei parametri per componente
    (esperimenti 02 e 03);
  \item \file{best\_model.pth} e \file{final\_model.pth} --- checkpoint
    selezionato sulla validation e stato dell'ultima epoca;
  \item \file{results.json} --- metriche, gap, storia per epoca ed eventuale
    robustezza al rumore;
  \item \file{events.out.tfevents.*} --- log TensorBoard.
\end{itemize}

\section{Controlli prima dell'analisi}

Prima dell'analisi si verificano completezza delle coppie, seed, configurazioni e
valori finiti nei \file{results.json}. Le tabelle derivano dagli aggregati JSON.
